\section{Discussion}
\label{sec:discussion}

Our data suggest three conditions for learning to add value over a strong classical controller, and only two of them were met. Under this protocol and in these settings, the conditions are, in order, as follows. First, headroom the classical toolbox cannot reach: it exists only on realistic whip (an oracle gap of +27.4\pp\ on TEST-C), while tracking ties and no juggling cell gained 10\pp\ from classical re-tuning in a small diagnostic. Second, a learner that passes the sanity gate on that gap at the control level: PPO and SAC did not on whip or juggling, and with the gate-chosen exploration scale PPO degraded S5 and S6 by 23.8 and 17.6\pp. Third, a planning step that does not trust one estimate. Robust plan selection, which scores candidates across several plausible plants, recovers 0.59 [0.36, 0.82] of the parameter gap with a classical Laplace ensemble and 0.90 [0.68, 1.16] with learned estimators, with the learned side ahead by +5.6\pp\ [0.6, 10.6]. The first and third conditions were met in our data; the second was not. Our reading is an inference rather than a test (Section~\ref{sec:k3}). We read the result as follows: the ingredient that helped was robust plan selection, a \emph{decision rule over a set of plausible plants} that either side can use, plus an advantage for learned estimators trained with simulator labels (+5.6\pp\ [0.6, 10.6]) that may reach 10\pp\ or more. We did not find a control-level learner (PPO, SAC, imitation student) that beats the strongest classical controller, and we have not shown that control-level learning is unattainable.

One piece of counter-evidence argues against reading K2 as ``classical is at least as good''. Under actuator noise four times larger, the classical T-adapt3-full lost 12.5\pp\ while the learned estimator did not: the paired difference is $-16.5$\pp\ [$-28.5$, $-5.0$] on HOLD ($n{=}200$, frozen methods, no re-tuning at the higher noise; Appendix~\ref{app:replication}). This is weak evidence, but together with the locked-set advantage of H16a it argues against a blanket reading of K2.

Three features could make the headline contrast of robust plan selection look cleaner than it is. The classical ensemble uses the generative parameter distribution as a prior, so it is not neutral; the learned estimators see the same distribution as training data. Both selection rules were picked on CAL among 20 or 22 variants, and the multi-start classical family, an obvious alternative, was degenerate in our implementation (it found no distinct second solution in 439 of 500 HOLD episodes). Finally, the success rate is close to chaotic in the estimated parameters, so differences of a few points between pipelines sit within the replicate noise we measured. The locked TEST-D set controls the selection effect but not the other two.

The null does not say that learning is worse. Three of eight settings had a learner that did not learn; the RL recipe was one family; 5 of the 8 settings have little or no headroom; and the +10\pp\ bar is ours. The comparison is also not equal on every axis: the learned estimators used 9--17$\times$ the classical episodes and simulator labels, whereas the classical ensemble costs about three times the CPU time per episode at test time.

For benchmark design, our results suggest four practices. Report the sanity gate and an oracle with every learned-versus-classical table, because a null from a learner that cannot learn a trivial correction is uninformative. Tag each result as pre-specified or post hoc; log locked-set accesses; and state which classical baseline is meant when the baseline changes during the study. Our own change of the strongest classical baseline on whip (R0 to C-ens) moved the learned advantage from +16.2\pp\ to +5.6\pp.
